\section{An invariant shows what cannot change}\label{ch05:sec:invariant}
Many problems describe a process: there is a collection of possible \emph{states}, and there are \emph{legal moves}, each of which turns one state into another. An \emph{invariant} of such a process is a quantity or property that every legal move leaves unchanged.

Here is an example. A row contains $n$ signs, each a plus or a minus, where $n\ge2$. The states of this process are the rows of $n$ signs. A legal move picks two different positions in the row and flips the sign in each: a plus becomes a minus, and a minus becomes a plus. Can we start from a row of all plus signs and reach a row with exactly one minus sign?

To find an invariant, replace each plus sign by the number $1$ and each minus sign by $-1$, and multiply all the numbers in the row together. Call the result the \emph{product of the row}. Flipping one sign multiplies that entry by $-1$, and therefore multiplies the whole product by $-1$. A legal move flips two signs, so it multiplies the product by $(-1)(-1)=1$: the product never changes. For instance, the row $+\,-\,+\,+$ has product $1\cdot(-1)\cdot1\cdot1=-1$. Flipping the first two signs gives $-\,+\,+\,+$, whose product is still $-1$.

A row of all plus signs has product $1$. Call a row \emph{reachable} if some sequence of legal moves turns the all-plus row into it. Induction on the number of moves shows that every reachable row also has product $1$. To set this up, let $P(k)$ be the statement ``every row obtained from the all-plus row by $k$ legal moves has product $1$.'' The base case $P(0)$ holds, because after $0$ moves the row is still all plus signs. For the induction step, assume $P(k)$ and take any row obtained by $k+1$ moves. It comes from a row obtained by $k$ moves, whose product is $1$ by the induction hypothesis, followed by one more move, which does not change the product. So its product is $1$, and $P(k+1)$ holds. A row with exactly one minus sign has product $-1$, so it can never be reached, however many moves we make.

Notice what this proof did not do: it never examined sequences of moves one by one. That would be hopeless, because moves can be repeated, so there are infinitely many sequences. Instead it found a quantity that no sequence of moves can change.

The product also has a simple meaning. If the row has $m$ minus signs, the plus signs contribute factors of $1$, so the product is $(-1)^m$. This is $1$ when $m$ is even and $-1$ when $m$ is odd. Seen this way, a legal move changes the number of minus signs by $+2$ (two plus signs flipped), by $-2$ (two minus signs flipped) or by $0$ (one of each). The number of minus signs starts at $0$, so it stays even.

Here the invariant is the only obstacle: every row with an even number of minus signs \emph{can} be reached from the all-plus row. Choose the positions that should hold minus signs, pair them up, and flip one pair at a time. In general, though, an invariant can only show that a state is unreachable; to show that a state can be reached, we must exhibit moves that actually get there.

An invariant on its own also does not show that a process comes to an end. Nothing in the sign-flipping rules ever forces us to stop: we can flip the first two signs, flip them back, and repeat forever, and every row along the way has product $1$. For a process that is meant to finish, we need a separate argument that it makes \emph{progress}. The next section studies such a process, Euclid's algorithm, and shows both kinds of argument working together.
